%COMANDO CHE DETERMINA IL TIPO DI DOCUMENTO CHE SI VUOLE CREARE

\documentclass[a4paper,12pt]{report}

%ELENCO DEI PACCHETTI UTILI PER LA STESURA DEL DOCUMENTO

\usepackage[utf8]{inputenc}
\usepackage[english, italian]{babel}
\usepackage{graphicx}
\usepackage{float}
\usepackage{tabularx}
\usepackage{makecell}
\usepackage{titlesec}
\usepackage{fancyhdr}
\usepackage{lastpage}
\usepackage{xurl}
\usepackage{hyperref}
\usepackage{geometry}
\usepackage[table,dvipsnames]{xcolor}
\setcounter{tocdepth}{5}
\setcounter{secnumdepth}{5}
\usepackage{caption}
\usepackage{etoolbox}% >= v2.1 2011-01-03
\usepackage{tikz}
\usepackage{listings}
\usepackage{scalerel}
\usepackage{longtable}
\usepackage{comment}
\usepackage{amssymb}
\usepackage{listings}

\titleformat{\chapter}[display]
{\normalfont\bfseries}{}{0pt}{\LARGE}

%COMANDO PER LA SPAZIATURA DEI TITOLI DAL BORDO DEL FOGLIO

\titlespacing*{\chapter}{0cm}{0cm}{0.2cm}

%COMANDO PER LA SPAZIATURA DEL TESTO DAI BORDI LATERALI

\geometry{
	left=20mm,
	right=20mm,
}

%COMADNO PER AVERE L'INDICE DEL NOME CHE SI VUOLE

\renewcommand{\contentsname}{Indice}

%COMADNI PER OTTENERE SUBSUBSECTION NUMERATE E PRESENTI NELL'INDICE

\setcounter{tocdepth}{5}
\setcounter{secnumdepth}{5}

%COMANDI PER OTTENERE HEADER E FOOTER

\pagestyle{plain}

\fancypagestyle{plain}{
	\fancyhf{}
	\lhead{\includegraphics[width=3cm]{../../immagini/minilogo.jpg}}
	\chead{}
	\rhead{\fontsize{12}{10}Verbale Esterno 2021-04-15}
	\lfoot{}
	\cfoot{\thepage\ di \pageref{LastPage}}
	\rfoot{}
}

\pagestyle{plain}

%COMANDI PER LINK

\hypersetup{
	colorlinks=true,
	linkcolor=black,
	filecolor=black,
	urlcolor=blue,
	citecolor=black,
}

%COMANDI PER SNIPPET DI CODICE

\BeforeBeginEnvironment{lstlisting}{\begin{mdframed}\vspace{-0.7em}}
	\AfterEndEnvironment{lstlisting}{\vspace{-0.5em}\end{mdframed}}

% needed for \lstcapt
\def\ifempty#1{\def\temparg{#1}\ifx\temparg\empty}

% make new caption command for listings
\usepackage{caption}
\newcommand{\lstcapt}[2][]{%
	\ifempty{#1}%
	\captionof{lstlisting}{#2}%
	\else%
	\captionof{lstlisting}[#1]{#2}%
	\fi%
	\vspace{0.75\baselineskip}%
}

\definecolor{atomlightorange}{rgb}{0.88,0.76,0.55}
\definecolor{atomdarkgrey}{RGB}{59,62,75}

% set listings
\lstset{%
	basicstyle=\footnotesize\ttfamily\color{atomlightorange},
	framesep=20pt,
	belowskip=10pt,
	aboveskip=10pt
}

% add frame environment
\usepackage[%
framemethod=tikz,
skipbelow=8pt,
skipabove=13pt
]{mdframed}
\mdfsetup{%
	leftmargin=0pt,
	rightmargin=0pt,
	backgroundcolor=atomdarkgrey,
	middlelinecolor=atomdarkgrey,
	roundcorner=6
}

%BIBLIOGRAFIA

\makeatletter
\def\thebibliography#1{\chapter*{Bibliografia\@mkboth
		{Bibliografia}{Bibliografia}}\list
	{[\arabic{enumi}]}{\settowidth\labelwidth{[#1]}\leftmargin\labelwidth
		\advance\leftmargin\labelsep
		\usecounter{enumi}}
	\def\newblock{\hskip .11em plus .33em minus .07em}
	\sloppy\clubpenalty4000\widowpenalty4000
	\sfcode`\.=1000\relax}
\makeatother

%DEFINIZIONE COLORI

\definecolor{airforceblue}{rgb}{0.36, 0.54, 0.66}


\begin{document}

\makeatletter
\begin{titlepage}
	\begin{center}
		\vspace*{-5cm}
		\author{Jawa Druids}
		\title{Nome documento}
		\date{} %LASCIARE QUESTO CAMPO VUOTO, SE LO TOLGO STAMPA LA DATA CORRENTE
		\includegraphics[width=0.7\linewidth]{../immagini/DRUIDSLOGO.jpg}\\[4ex]
		{\huge \bfseries  \@title }\\[2ex]
		{\LARGE  \@author}\\[50ex]
		\vspace*{-9cm}
		\begin{table}[H]
			\renewcommand{\arraystretch}{1.4}
			\centering
			\begin{tabular}{r | l}
				\textbf{Versione} & x.x.x \\%RIGA PER INSERIRE LA VERSIONE ULTIMA DEL DOCUMENTO
				\textbf{Data approvazione} & xx-xx-xxxx\\
				\textbf{Responsabile} & Nome Cognome\\
				\textbf{Redattori} & Andrea Dorigo \\
				\textbf{Verificatori} & Nome Cognome \\
		%MAKECELL SERVE PER POI ANDARE A CAPO ALL'INTERNO DELLA CELLA
				\textbf{Stato} & Redazione in corso\\
				\textbf{Lista distribuzione} & \makecell[tl]{Jawa Druids \\ prof. Tullio Vardanega \\ prof. Riccardo Cardin \\ Sync Lab}\\
				\textbf{Uso} & Esterno
			\end{tabular}
		\end{table}
		\vspace{0.1cm}
		\hfill \break
		\fontsize{17}{10}\textbf{Sommario} \\
		\vspace{0.1cm}
    Il presente documento contiene la pianificazione delle attività del gruppo Jawa Druids atte al soddisfacimento del capitolato \normalsize\textit{GDP: Gathering Detection Platform} di Sync Lab.
	\end{center}
\end{titlepage}
\makeatother

\quad
\begin{center}
	\LARGE\textbf{Registro delle modifiche}
\end{center}

\def\tabularxcolumn#1{m{#1}}
{\rowcolors{2}{RawSienna!90!RawSienna!20}{RawSienna!70!RawSienna!40}


\begin{center}
	\renewcommand{\arraystretch}{1.4}
	\begin{longtable}[c]{|p{2cm-1\tabcolsep}|p{2cm}|p{3cm-2\tabcolsep}|p{2,5cm-2\tabcolsep}|p{2,5cm}|p{4cm-2\tabcolsep}|}
		\hline
		\rowcolor{airforceblue}
		\makecell[c]{\textbf{Versione}} & \makecell[c]{\textbf{Data}} & \makecell[c]{\textbf{Autore}} & \makecell[c]{\textbf{Ruolo}} & \makecell[c]{\textbf{Verificatore}} & \makecell[c]{\textbf{Modifica}}\\
		\hline
		\centering v1.0.0 & 07-02-2021 & Andrea Dorigo & \centering \textit{Analista} & \centering - & \textit{Approvazione del documento} \\
		\hline
		\centering v0.1.0 & 05-02-2021 & Emma Roveroni & \centering \textit{Analista} & Andrea Cecchin & \textit{Revisione complessiva del documento} \\
		\hline
		\centering v0.0.1 & 03-02-2021 & Emma Roveroni & \centering \textit{Analista} & Andrea Cecchin & \textit{Stesura del documento} \\
		\hline
		
	\end{longtable}
\end{center}

%COMANDO PER LA CREAZIONE DELL'INDICE

\tableofcontents{}
\listoffigures
\listoftables
\chapter{Introduzione}

Lo scopo di questo documento è la stesura di un elenco di linee guida ed esempi che i componenti sono incoraggiati a seguire per migliorare l'efficacia della collaborazione.
A differenza delle Norme di Progetto, questo documento \`{e} redatto in un linguaggio pi\`{u} informale per facilitarne la comprensione, con spezzati di codice a scopo esemplificativo e riferimenti esterni sulle best practices da seguire.\\
Il documento pu\`{o} essere soggetto a modifiche ed aggiunte per tutta la durata del progetto.

\chapter{Analisi dei rischi}\label{AnalisiDeiRischi}

\section{Piano per la gestione dei rischi}\label{AnalisiDeiRischiPianoPerLaGestioneDeiRischi}
Con l'intento di prevenire il naturale insorgere di problemi durante lo svolgimento del progetto è stato elaborato un'approfondito piano per la gestione dei rischi. Quest'ultimo è suddiviso in quattro attività$_{\scaleto{G}{3pt}}$:
\begin{itemize}
  \item \textbf{Individuazione dei rischi:} attività$_{\scaleto{G}{3pt}}$ di identificazione e documentazione di possibili elementi problematici che possano ostacolare il naturale percorso del progetto;
  \item \textbf{Analisi dei rischi:} attività$_{\scaleto{G}{3pt}}$ di analisi dei fattori di rischio, che si articola in probabilità di occorrenza, indice di gravità e conseguente impatto sul progetto;
  \item \textbf{Pianificazione di controllo:} attività$_{\scaleto{G}{3pt}}$ di pianificazione delle misure da adottare per la prevenzione e contenimento del problema;
  \item \textbf{Monitoraggio dei rischi:} attività$_{\scaleto{G}{3pt}}$ di controllo dei rischi che accompagna tutto lo svolgimento del progetto, al fine di evitarli o agire tempestivamente alla loro occorrenza per contenerne i danni.
\end{itemize}
Le principali tipologie di rischio sono state quindi codificate e categorizzate come segue:
\begin{itemize}
  \item \textbf{RT:} Rischi legati alle tecnologie;
  \item \textbf{RO:} Rischi legati all'organizzazione;
  \item \textbf{RI:} Rischi interpersonali, ovvero legati alle relazioni personali interne ed esterne o alla disponibilità e risorse dei componenti.
\end{itemize}

\quad
\begin{center}
	\LARGE\textbf{Rischi legati alle tecnologie}
\end{center}

\def\tabularxcolumn#1{m{#1}}
{\rowcolors{2}{RawSienna!90!RawSienna!20}{RawSienna!70!RawSienna!40}

	\begin{center}
		\renewcommand{\arraystretch}{1.4}
		\begin{tabularx}{\textwidth}{|c|X|}
			\hline
			\rowcolor{airforceblue}
			\multicolumn{2}{|c|}{\textit{Inesperienza tecnologica}}\\
			\hline
			\textit{Codice} & RT1 \\
			\hline
			\textit{Descrizione} & Alcune tecnologie utilizzate in questo progetto sono nuove per tutti i membri del gruppo di lavoro. \\
			\hline
			\textit{Conseguenza} & Lo studio e l'apprendimento di tali tecnologie potrebbero richiedere un intervallo di tempo difficile da quantificare, maggiore del previsto e variabile da membro a membro con conseguenti difficoltà operative. \\
			\hline
			\textit{Possibilità di occorrenza} & Alta. \\
			\hline
			\textit{Pericolosità} & Alta. \\
			\hline
			\textit{Precauzioni} & Il \textit{Responsabile di Progetto} dovrà suddividere i compiti nel modo più congruo possibile, considerando le conoscenze preliminari di ciascun componente; prevederà inoltre un tempo di Slack$_G$ maggiore per i compiti assegnati ad un componente senza particolare famigliarità con la relativa tecnologia. Il \textit{Responsabile di Progetto} assegnerà i task$_G$ di maggiore complessità a più membri ove necessario.  \\
			\hline
			\textit{Piano di contingenza} & Ciascun membro comunicherà il prima possibile al \textit{Responsabile di progetto} la previsione di un eventuale ritardo o mancanza; egli provvederà a ridistribuire i compiti se necessario in modo da sanare eventuali lacune o sottostime. \\
			\hline
		\end{tabularx}
	\captionof{table}{\textbf{Analisi dei rischi delle tecnologie utilizzate}}
	\end{center}


\def\tabularxcolumn#1{m{#1}}
{\rowcolors{2}{RawSienna!90!RawSienna!20}{RawSienna!70!RawSienna!40}

	\begin{center}
		\renewcommand{\arraystretch}{1.4}
		\begin{tabularx}{\textwidth}{|c|X|}
			\hline
			\rowcolor{airforceblue}
			\multicolumn{2}{|c|}{\textit{Software terze parti}}\\
			\hline
			\textit{Codice} & RT2 \\
			\hline
			\textit{Descrizione} & Eventuali problematiche con software di terze parti, quali la mancanza di documentazione o problemi tecnici, sono indipendenti dai membri del gruppo. \\
			\hline
			\textit{Conseguenza} & Ciò causerebbe ritardi pesanti sul proseguo del lavoro e anche possibili ritardi sulla consegna.
			Necessità di cambiare tecnologia, potrebbe richiedere molto tempo e risorse per la ricerca di una sostituzione. \\
			\hline
			\textit{Possibilità di occorrenza} & Bassa. \\
			\hline
			\textit{Pericolosità} & Alta. \\
			\hline
			\textit{Precauzioni} & Il gruppo sceglierà i software più stabili e documentati per evitare questi tipi di problemi.  \\
			\hline
			\textit{Piano di contingenza} & Assieme al \textit{Responsabile di progetto} il gruppo di lavoro si attiverà al fine di tentare di risolvere il problema. Se ciò non è possibile sarà necessario un cambio di tecnologia anche tramite l'aiuto del proponente$_{\scaleto{G}{3pt}}$.  \\
			\hline
		\end{tabularx}
		\captionof{table}{\textbf{Analisi dei rischi dei software di terze parti}}
	\end{center}


\def\tabularxcolumn#1{m{#1}}
{\rowcolors{2}{RawSienna!90!RawSienna!20}{RawSienna!70!RawSienna!40}

	\begin{center}
		\renewcommand{\arraystretch}{1.4}
		\begin{tabularx}{\textwidth}{|c|X|}
			\hline
			\rowcolor{airforceblue}
			\multicolumn{2}{|c|}{\textit{Validità dei dati}}\\
			\hline
			\textit{Codice} & RT3 \\
			\hline
			\textit{Descrizione} & Problemi legati alla validità e all'elaborazione dei dati. \\
			\hline
			\textit{Conseguenza} & Arresto obbligato del lavoro in corso, con possibilità di invalidazione del lavoro svolto fino a quel momento. \\
			\hline
			\textit{Possibilità di occorrenza} & Medio/Alta. \\
			\hline
			\textit{Pericolosità} & Molto alta. \\
			\hline
			\textit{Precauzioni} & Prima dell'inizio della raccolta dati il gruppo si assicurerà che la fonte sia affidabile e coerente.
			Questa operazione sarà svolta per prima in quanto critica all'intero sviluppo.  \\
			\hline
			\textit{Piano di contingenza} & Il gruppo, insieme al proponente$_{\scaleto{G}{3pt}}$, valuterà se sarà necessario cambiare solo la fonte di provenienza dei dati oppure simularli in maniera consona. \\
			\hline
		\end{tabularx}
		\captionof{table}{\textbf{Analisi dei rischi della validità dei dati}}
	\end{center}

\quad
\begin{center}
	\LARGE\textbf{Rischi legati all'organizzazione}
\end{center}

\def\tabularxcolumn#1{m{#1}}
{\rowcolors{2}{RawSienna!90!RawSienna!20}{RawSienna!70!RawSienna!40}

	\begin{center}
		\renewcommand{\arraystretch}{1.4}
		\begin{tabularx}{\textwidth}{|c|X|}
			\hline
			\rowcolor{airforceblue}
			\multicolumn{2}{|c|}{\textit{Problemi organizzativi}}\\
			\hline
			\textit{Codice} & RO1 \\
			\hline
			\textit{Descrizione} & I problemi organizzativi possono scaturire da vari motivi, sia da parte dei membri che dal proponente$_{\scaleto{G}{3pt}}$, così come dagli impegni personali e dai periodi vacanzieri.  \\
			\hline
			\textit{Conseguenza} & Questi problemi possono far ritardare il completamento dei tasks$_{\scaleto{G}{3pt}}$ di un tempo più o meno definito, rendendo l'avanzamento più lento o addirittura bloccandolo. \\
			\hline
			\textit{Possibilità di occorrenza} & Alta. \\
			\hline
			\textit{Pericolosità} & Alta. \\
			\hline
			\textit{Precauzioni} & Ogni membro del gruppo di lavoro dovrà avvisare il \textit{Responsabile di progetto} nel caso in cui, per cause di forza maggiore, non si riesca a completare il task$_{\scaleto{G}{3pt}}$ nel tempo deciso oppure non si riesca proprio a farlo. \\
			\hline
			\textit{Piano di contingenza} & Il \textit{Responsabile di progetto} avrà l'incarico di riassegnare i compiti in modo da riuscire a completare i task$_{\scaleto{G}{3pt}}$ nel tempo stimato, così da non avere ritardi nel portarli a termine.
			Nel caso in cui sia il proponente$_{\scaleto{G}{3pt}}$ a creare questi disagi organizzativi, sarà sempre premura del \textit{Responsabile di progetto} risolvere il problema mediante i canali di comunicazione adatti.  \\
			\hline
		\end{tabularx}
	\captionof{table}{\textbf{Analisi dei rischi dei problemi organizzativi}}
	\end{center}


\def\tabularxcolumn#1{m{#1}}
{\rowcolors{2}{RawSienna!90!RawSienna!20}{RawSienna!70!RawSienna!40}

	\begin{center}
		\renewcommand{\arraystretch}{1.4}
		\begin{tabularx}{\textwidth}{|c|X|}
			\hline
			\rowcolor{airforceblue}
			\multicolumn{2}{|c|}{\textit{Problemi dei sistemi operativi e configurazioni software}}\\
			\hline
			\textit{Codice} & RO2 \\
			\hline
			\textit{Descrizione} & Problemi che possono sorgere per via delle differenze degli standard utilizzati dai software in base al sistema operativo in cui sono installati. \\
			\hline
			\textit{Conseguenza} & Possibili incongruenze nelle funzionalità o visualizzazione del prodotto software.  \\
			\hline
			\textit{Possibilità di occorrenza} & Media. \\
			\hline
			\textit{Pericolosità} & Media. \\
			\hline
			\textit{Precauzioni} & Il gruppo cercherà di trovare una configurazione software adatta per ogni sistema operativo in modo da ridurre al minimo il pericolo. \\
			\hline
			\textit{Piano di contingenza} & Il gruppo cercherà di trovare una soluzione nel minor tempo possibile. \\
			\hline
		\end{tabularx}
		\captionof{table}{\textbf{Analisi dei rischi su software e sistemi operativi}}
	\end{center}


\quad
\begin{center}
	\LARGE\textbf{Rischi interpersonali}
\end{center}

\def\tabularxcolumn#1{m{#1}}
{\rowcolors{2}{RawSienna!90!RawSienna!20}{RawSienna!70!RawSienna!40}

	\begin{center}
		\renewcommand{\arraystretch}{1.4}

		\begin{tabularx}{\textwidth}{|c|X|}
			\hline
			\rowcolor{airforceblue}
			\multicolumn{2}{|c|}{\textit{Problemi di relazione tra i membri}}\\
			\hline
			\textit{Codice} & RI1 \\
			\hline
			\textit{Descrizione} & Problemi legati ai contrasti che potrebbero intercorrere tra i membri. \\
			\hline
			\textit{Conseguenza} & Difficoltà di avanzamento nel lavoro, poca collaborazione tra i membri in contrasto, malumore nel gruppo. \\
			\hline
			\textit{Possibilità di occorrenza} & Media \\
			\hline
			\textit{Pericolosità} & Alta. \\
			\hline
			\textit{Precauzioni} & In caso di nascite di contrasti tra i membri, questi dovranno immediatamente coinvolgere il \textit{Responsabile di progetto} in modo da poter risolvere subito la diatriba.
			In caso non riesca a risolvere la controversia,  comunicherà col \textit{prof. Tullio Vardanega} per la risoluzione dei problemi. \\
			\hline
			\textit{Piano di contingenza} & I membri dovranno impegnarsi nel ridurre al minimo eventuali tensioni tra di loro per favorire l'avanzamento dei lavori e per realizzare al meglio il progetto.\\
			\hline
		\end{tabularx}
	\captionof{table}{\textbf{Analisi dei rischi dei problemi relazionali}}
	\end{center}


%PER RENDERE PIÙ CHIARA LA STESURA DEI DOCUMENTI È MEGLIO LASCIARE SEPARATI IN FILE DIVERSI OGNI CAPITOLO

% \input{esempio} -- esempio di codice per inserire un nuovo capitolo

\end{document}

% \documentclass[a4paper,12pt]{report}
% \usepackage[utf8]{inputenc}
% \usepackage{graphicx}
% \usepackage{float}
% \usepackage{tabularx}
% \usepackage{makecell}
% \usepackage{titlesec}
% \usepackage{fancyhdr}
% \usepackage{lastpage}
% \usepackage{xurl}
% \usepackage{hyperref}
% \usepackage{geometry}
% \usepackage{color}
% \usepackage{microtype}
% \usepackage{enumerate}
% \usepackage{listings}
% \usepackage{tabularx}
% \usepackage[table,dvipsnames]{xcolor}
% \usepackage{caption}
% \captionsetup{labelformat=empty}
% \setcounter{tocdepth}{5}
% \setcounter{secnumdepth}{5} %SONO I DUE COMANDI PER AVERE LE SUBSUBSECTION NUMERATE E PRESENTI NELL'INDICE
% \renewcommand{\contentsname}{Indice} %QUESTO SERVE PER AVERE L'INDICE CON IL NOME CHE VOGLIO IO
% \renewcommand{\listfigurename}{Lista delle figure}
% \renewcommand{\listtablename}{Lista delle tabelle}
% \titleformat{\chapter}[display]
% {\normalfont\bfseries}{}{0pt}{\LARGE} %QUESTO SERVE PER AVERE SOLO IL NOME DEL CAPITOLO CHE VOGLIO IO
% \titlespacing*{\chapter}{0cm}{0cm}{0.2cm}
% \geometry{
% 	left=20mm,
% 	right=20mm,
% }
% \fancypagestyle{plain}{
% 	\fancyhf{}
% 	\lhead{\includegraphics[width=3cm]{../immagini/minilogo.jpg}}
% 	\chead{}
% 	\rhead{\fontsize{12}{10}Piano di progetto}
% 	\lfoot{}
% 	\cfoot{\thepage\ di \pageref*{LastPage}}
% 	\rfoot{}
% }
%
% \definecolor{atomlightorange}{rgb}{0.88,0.76,0.55}
% \definecolor{atomdarkgrey}{RGB}{59,62,75}
%
% \usepackage{tikz}
%
% % set listings
% \lstset{%
%     basicstyle=\footnotesize\ttfamily\color{atomlightorange},
%     framesep=20pt,
% 		belowskip=10pt,
% 		aboveskip=10pt
% }
%
% % add frame environment
% \usepackage[%
%     framemethod=tikz,
%     skipbelow=8pt,
%     skipabove=13pt
% ]{mdframed}
% \mdfsetup{%
%     leftmargin=0pt,
%     rightmargin=0pt,
%     backgroundcolor=atomdarkgrey,
%     middlelinecolor=atomdarkgrey,
%     roundcorner=6
% }
%
% \usepackage{etoolbox}% >= v2.1 2011-01-03
% \BeforeBeginEnvironment{lstlisting}{\begin{mdframed}\vspace{-0.7em}}
% \AfterEndEnvironment{lstlisting}{\vspace{-0.5em}\end{mdframed}}
%
% % needed for \lstcapt
% \def\ifempty#1{\def\temparg{#1}\ifx\temparg\empty}
%
% % make new caption command for listings
% \usepackage{caption}
% \newcommand{\lstcapt}[2][]{%
%     \ifempty{#1}%
%         \captionof{lstlisting}{#2}%
%     \else%
%         \captionof{lstlisting}[#1]{#2}%
%     \fi%
%     \vspace{0.75\baselineskip}%
% }
%
% \hypersetup{
%     colorlinks=true,
%     linkcolor=black,
%     filecolor=black,
%     urlcolor=blue,
% 		citecolor=black,
% }
%
% \pagestyle{plain}
%
% \makeatletter
% 	 \def\thebibliography#1{\chapter*{Bibliografia\@mkboth
% 		 {Bibliografia}{Bibliografia}}\list
% 		 {[\arabic{enumi}]}{\settowidth\labelwidth{[#1]}\leftmargin\labelwidth
%  \advance\leftmargin\labelsep
%  \usecounter{enumi}}
%  \def\newblock{\hskip .11em plus .33em minus .07em}
%  \sloppy\clubpenalty4000\widowpenalty4000
%  \sfcode`\.=1000\relax}
% 	 \makeatother
%
% \newcolumntype{Y}{>{\centering\arraybackslash}X}
%
% \begin{document}
%
% \makeatletter
% \begin{titlepage}
%     \begin{center}
%     \vspace*{-5,0cm}
%     \author{Jawa Druids}
%     \title{Piano di progetto}
%     \date{} %LASCIARE QUESTO CAMPO VUOTO, SE LO TOLGO STAMPA LA DATA CORRENTE
%     \includegraphics[width=0.7\linewidth]{../immagini/DRUIDSLOGO.jpg}\\[4ex]
%     {\huge \bfseries  \@title }\\[2ex]
%     {\LARGE  \@author}\\[50ex]
%     \vspace*{-8,0cm}
%     \begin{table}[H]
%         \centering
%         \begin{tabular}{c|c}
%             \textbf{Versione} & x.x.x \\
%             \textbf{Data approvazione} & xx-xx-xxxx\\
%             \textbf{Responsabile} & Nome Cognome\\
%             \textbf{Redattori} & Andrea Dorigo \\
%             \textbf{Verificatori} & Nome Cognome \\
%         %MAKECELL SERVE PER POI ANDARE A CAPO ALL'INTERNO DELLA CELLA
%             \textbf{Stato} & Redazione in corso\\
%             \textbf{Lista distribuzione} & \makecell{Jawa Druids \\ prof. Tullio Vardanega \\ prof. Riccardo Cardin \\ Sync Lab}\\
%             \textbf{Uso} & Esterno
%         \end{tabular}
%     \end{table}
% 					\vspace{.4cm}
% 	\hfill \break
%     \fontsize{17}{10}\textbf{Sommario}\\
% 		\vspace{.3cm}
%     Il presente documento contiene la pianificazione delle attività del gruppo Jawa Druids atte al soddisfacimento del capitolato \normalsize\textit{GDP: Gathering Detection Platform} di Sync Lab.
%     \end{center}
% \end{titlepage}
% \makeatother
%
% \quad
\begin{center}
	\LARGE\textbf{Registro delle modifiche}
\end{center}

\def\tabularxcolumn#1{m{#1}}
{\rowcolors{2}{RawSienna!90!RawSienna!20}{RawSienna!70!RawSienna!40}


\begin{center}
	\renewcommand{\arraystretch}{1.4}
	\begin{longtable}[c]{|p{2cm-1\tabcolsep}|p{2cm}|p{3cm-2\tabcolsep}|p{2,5cm-2\tabcolsep}|p{2,5cm}|p{4cm-2\tabcolsep}|}
		\hline
		\rowcolor{airforceblue}
		\makecell[c]{\textbf{Versione}} & \makecell[c]{\textbf{Data}} & \makecell[c]{\textbf{Autore}} & \makecell[c]{\textbf{Ruolo}} & \makecell[c]{\textbf{Verificatore}} & \makecell[c]{\textbf{Modifica}}\\
		\hline
		\centering v1.0.0 & 07-02-2021 & Andrea Dorigo & \centering \textit{Analista} & \centering - & \textit{Approvazione del documento} \\
		\hline
		\centering v0.1.0 & 05-02-2021 & Emma Roveroni & \centering \textit{Analista} & Andrea Cecchin & \textit{Revisione complessiva del documento} \\
		\hline
		\centering v0.0.1 & 03-02-2021 & Emma Roveroni & \centering \textit{Analista} & Andrea Cecchin & \textit{Stesura del documento} \\
		\hline
		
	\end{longtable}
\end{center}
% \tableofcontents
% \listoffigures
% \listoftables
% \chapter{Introduzione}

Lo scopo di questo documento è la stesura di un elenco di linee guida ed esempi che i componenti sono incoraggiati a seguire per migliorare l'efficacia della collaborazione.
A differenza delle Norme di Progetto, questo documento \`{e} redatto in un linguaggio pi\`{u} informale per facilitarne la comprensione, con spezzati di codice a scopo esemplificativo e riferimenti esterni sulle best practices da seguire.\\
Il documento pu\`{o} essere soggetto a modifiche ed aggiunte per tutta la durata del progetto.

% \chapter{Analisi dei rischi}\label{AnalisiDeiRischi}

\section{Piano per la gestione dei rischi}\label{AnalisiDeiRischiPianoPerLaGestioneDeiRischi}
Con l'intento di prevenire il naturale insorgere di problemi durante lo svolgimento del progetto è stato elaborato un'approfondito piano per la gestione dei rischi. Quest'ultimo è suddiviso in quattro attività$_{\scaleto{G}{3pt}}$:
\begin{itemize}
  \item \textbf{Individuazione dei rischi:} attività$_{\scaleto{G}{3pt}}$ di identificazione e documentazione di possibili elementi problematici che possano ostacolare il naturale percorso del progetto;
  \item \textbf{Analisi dei rischi:} attività$_{\scaleto{G}{3pt}}$ di analisi dei fattori di rischio, che si articola in probabilità di occorrenza, indice di gravità e conseguente impatto sul progetto;
  \item \textbf{Pianificazione di controllo:} attività$_{\scaleto{G}{3pt}}$ di pianificazione delle misure da adottare per la prevenzione e contenimento del problema;
  \item \textbf{Monitoraggio dei rischi:} attività$_{\scaleto{G}{3pt}}$ di controllo dei rischi che accompagna tutto lo svolgimento del progetto, al fine di evitarli o agire tempestivamente alla loro occorrenza per contenerne i danni.
\end{itemize}
Le principali tipologie di rischio sono state quindi codificate e categorizzate come segue:
\begin{itemize}
  \item \textbf{RT:} Rischi legati alle tecnologie;
  \item \textbf{RO:} Rischi legati all'organizzazione;
  \item \textbf{RI:} Rischi interpersonali, ovvero legati alle relazioni personali interne ed esterne o alla disponibilità e risorse dei componenti.
\end{itemize}

\quad
\begin{center}
	\LARGE\textbf{Rischi legati alle tecnologie}
\end{center}

\def\tabularxcolumn#1{m{#1}}
{\rowcolors{2}{RawSienna!90!RawSienna!20}{RawSienna!70!RawSienna!40}

	\begin{center}
		\renewcommand{\arraystretch}{1.4}
		\begin{tabularx}{\textwidth}{|c|X|}
			\hline
			\rowcolor{airforceblue}
			\multicolumn{2}{|c|}{\textit{Inesperienza tecnologica}}\\
			\hline
			\textit{Codice} & RT1 \\
			\hline
			\textit{Descrizione} & Alcune tecnologie utilizzate in questo progetto sono nuove per tutti i membri del gruppo di lavoro. \\
			\hline
			\textit{Conseguenza} & Lo studio e l'apprendimento di tali tecnologie potrebbero richiedere un intervallo di tempo difficile da quantificare, maggiore del previsto e variabile da membro a membro con conseguenti difficoltà operative. \\
			\hline
			\textit{Possibilità di occorrenza} & Alta. \\
			\hline
			\textit{Pericolosità} & Alta. \\
			\hline
			\textit{Precauzioni} & Il \textit{Responsabile di Progetto} dovrà suddividere i compiti nel modo più congruo possibile, considerando le conoscenze preliminari di ciascun componente; prevederà inoltre un tempo di Slack$_G$ maggiore per i compiti assegnati ad un componente senza particolare famigliarità con la relativa tecnologia. Il \textit{Responsabile di Progetto} assegnerà i task$_G$ di maggiore complessità a più membri ove necessario.  \\
			\hline
			\textit{Piano di contingenza} & Ciascun membro comunicherà il prima possibile al \textit{Responsabile di progetto} la previsione di un eventuale ritardo o mancanza; egli provvederà a ridistribuire i compiti se necessario in modo da sanare eventuali lacune o sottostime. \\
			\hline
		\end{tabularx}
	\captionof{table}{\textbf{Analisi dei rischi delle tecnologie utilizzate}}
	\end{center}


\def\tabularxcolumn#1{m{#1}}
{\rowcolors{2}{RawSienna!90!RawSienna!20}{RawSienna!70!RawSienna!40}

	\begin{center}
		\renewcommand{\arraystretch}{1.4}
		\begin{tabularx}{\textwidth}{|c|X|}
			\hline
			\rowcolor{airforceblue}
			\multicolumn{2}{|c|}{\textit{Software terze parti}}\\
			\hline
			\textit{Codice} & RT2 \\
			\hline
			\textit{Descrizione} & Eventuali problematiche con software di terze parti, quali la mancanza di documentazione o problemi tecnici, sono indipendenti dai membri del gruppo. \\
			\hline
			\textit{Conseguenza} & Ciò causerebbe ritardi pesanti sul proseguo del lavoro e anche possibili ritardi sulla consegna.
			Necessità di cambiare tecnologia, potrebbe richiedere molto tempo e risorse per la ricerca di una sostituzione. \\
			\hline
			\textit{Possibilità di occorrenza} & Bassa. \\
			\hline
			\textit{Pericolosità} & Alta. \\
			\hline
			\textit{Precauzioni} & Il gruppo sceglierà i software più stabili e documentati per evitare questi tipi di problemi.  \\
			\hline
			\textit{Piano di contingenza} & Assieme al \textit{Responsabile di progetto} il gruppo di lavoro si attiverà al fine di tentare di risolvere il problema. Se ciò non è possibile sarà necessario un cambio di tecnologia anche tramite l'aiuto del proponente$_{\scaleto{G}{3pt}}$.  \\
			\hline
		\end{tabularx}
		\captionof{table}{\textbf{Analisi dei rischi dei software di terze parti}}
	\end{center}


\def\tabularxcolumn#1{m{#1}}
{\rowcolors{2}{RawSienna!90!RawSienna!20}{RawSienna!70!RawSienna!40}

	\begin{center}
		\renewcommand{\arraystretch}{1.4}
		\begin{tabularx}{\textwidth}{|c|X|}
			\hline
			\rowcolor{airforceblue}
			\multicolumn{2}{|c|}{\textit{Validità dei dati}}\\
			\hline
			\textit{Codice} & RT3 \\
			\hline
			\textit{Descrizione} & Problemi legati alla validità e all'elaborazione dei dati. \\
			\hline
			\textit{Conseguenza} & Arresto obbligato del lavoro in corso, con possibilità di invalidazione del lavoro svolto fino a quel momento. \\
			\hline
			\textit{Possibilità di occorrenza} & Medio/Alta. \\
			\hline
			\textit{Pericolosità} & Molto alta. \\
			\hline
			\textit{Precauzioni} & Prima dell'inizio della raccolta dati il gruppo si assicurerà che la fonte sia affidabile e coerente.
			Questa operazione sarà svolta per prima in quanto critica all'intero sviluppo.  \\
			\hline
			\textit{Piano di contingenza} & Il gruppo, insieme al proponente$_{\scaleto{G}{3pt}}$, valuterà se sarà necessario cambiare solo la fonte di provenienza dei dati oppure simularli in maniera consona. \\
			\hline
		\end{tabularx}
		\captionof{table}{\textbf{Analisi dei rischi della validità dei dati}}
	\end{center}

\quad
\begin{center}
	\LARGE\textbf{Rischi legati all'organizzazione}
\end{center}

\def\tabularxcolumn#1{m{#1}}
{\rowcolors{2}{RawSienna!90!RawSienna!20}{RawSienna!70!RawSienna!40}

	\begin{center}
		\renewcommand{\arraystretch}{1.4}
		\begin{tabularx}{\textwidth}{|c|X|}
			\hline
			\rowcolor{airforceblue}
			\multicolumn{2}{|c|}{\textit{Problemi organizzativi}}\\
			\hline
			\textit{Codice} & RO1 \\
			\hline
			\textit{Descrizione} & I problemi organizzativi possono scaturire da vari motivi, sia da parte dei membri che dal proponente$_{\scaleto{G}{3pt}}$, così come dagli impegni personali e dai periodi vacanzieri.  \\
			\hline
			\textit{Conseguenza} & Questi problemi possono far ritardare il completamento dei tasks$_{\scaleto{G}{3pt}}$ di un tempo più o meno definito, rendendo l'avanzamento più lento o addirittura bloccandolo. \\
			\hline
			\textit{Possibilità di occorrenza} & Alta. \\
			\hline
			\textit{Pericolosità} & Alta. \\
			\hline
			\textit{Precauzioni} & Ogni membro del gruppo di lavoro dovrà avvisare il \textit{Responsabile di progetto} nel caso in cui, per cause di forza maggiore, non si riesca a completare il task$_{\scaleto{G}{3pt}}$ nel tempo deciso oppure non si riesca proprio a farlo. \\
			\hline
			\textit{Piano di contingenza} & Il \textit{Responsabile di progetto} avrà l'incarico di riassegnare i compiti in modo da riuscire a completare i task$_{\scaleto{G}{3pt}}$ nel tempo stimato, così da non avere ritardi nel portarli a termine.
			Nel caso in cui sia il proponente$_{\scaleto{G}{3pt}}$ a creare questi disagi organizzativi, sarà sempre premura del \textit{Responsabile di progetto} risolvere il problema mediante i canali di comunicazione adatti.  \\
			\hline
		\end{tabularx}
	\captionof{table}{\textbf{Analisi dei rischi dei problemi organizzativi}}
	\end{center}


\def\tabularxcolumn#1{m{#1}}
{\rowcolors{2}{RawSienna!90!RawSienna!20}{RawSienna!70!RawSienna!40}

	\begin{center}
		\renewcommand{\arraystretch}{1.4}
		\begin{tabularx}{\textwidth}{|c|X|}
			\hline
			\rowcolor{airforceblue}
			\multicolumn{2}{|c|}{\textit{Problemi dei sistemi operativi e configurazioni software}}\\
			\hline
			\textit{Codice} & RO2 \\
			\hline
			\textit{Descrizione} & Problemi che possono sorgere per via delle differenze degli standard utilizzati dai software in base al sistema operativo in cui sono installati. \\
			\hline
			\textit{Conseguenza} & Possibili incongruenze nelle funzionalità o visualizzazione del prodotto software.  \\
			\hline
			\textit{Possibilità di occorrenza} & Media. \\
			\hline
			\textit{Pericolosità} & Media. \\
			\hline
			\textit{Precauzioni} & Il gruppo cercherà di trovare una configurazione software adatta per ogni sistema operativo in modo da ridurre al minimo il pericolo. \\
			\hline
			\textit{Piano di contingenza} & Il gruppo cercherà di trovare una soluzione nel minor tempo possibile. \\
			\hline
		\end{tabularx}
		\captionof{table}{\textbf{Analisi dei rischi su software e sistemi operativi}}
	\end{center}


\quad
\begin{center}
	\LARGE\textbf{Rischi interpersonali}
\end{center}

\def\tabularxcolumn#1{m{#1}}
{\rowcolors{2}{RawSienna!90!RawSienna!20}{RawSienna!70!RawSienna!40}

	\begin{center}
		\renewcommand{\arraystretch}{1.4}

		\begin{tabularx}{\textwidth}{|c|X|}
			\hline
			\rowcolor{airforceblue}
			\multicolumn{2}{|c|}{\textit{Problemi di relazione tra i membri}}\\
			\hline
			\textit{Codice} & RI1 \\
			\hline
			\textit{Descrizione} & Problemi legati ai contrasti che potrebbero intercorrere tra i membri. \\
			\hline
			\textit{Conseguenza} & Difficoltà di avanzamento nel lavoro, poca collaborazione tra i membri in contrasto, malumore nel gruppo. \\
			\hline
			\textit{Possibilità di occorrenza} & Media \\
			\hline
			\textit{Pericolosità} & Alta. \\
			\hline
			\textit{Precauzioni} & In caso di nascite di contrasti tra i membri, questi dovranno immediatamente coinvolgere il \textit{Responsabile di progetto} in modo da poter risolvere subito la diatriba.
			In caso non riesca a risolvere la controversia,  comunicherà col \textit{prof. Tullio Vardanega} per la risoluzione dei problemi. \\
			\hline
			\textit{Piano di contingenza} & I membri dovranno impegnarsi nel ridurre al minimo eventuali tensioni tra di loro per favorire l'avanzamento dei lavori e per realizzare al meglio il progetto.\\
			\hline
		\end{tabularx}
	\captionof{table}{\textbf{Analisi dei rischi dei problemi relazionali}}
	\end{center}

% % \begin{thebibliography}{}
    % \bibitem{GitFlowCheatseet}
    % Daniel Kummer.
    % \textit{Git flow cheatsheet - efficient branching using git-flow by Vincent Driessen}.\\
		% \url{https://danielkummer.github.io/git-flow-cheatsheet/}.
		% \bibitem{GitFlowWorkflowTutorial}
		% Atlassian Bitbucket.
		% \textit{Git flow workflow tutorial}.\\
		% \url{https://www.atlassian.com/git/tutorials/comparing-workflows/gitflow-workflow}.
		% \bibitem{GitCommitGuidelines}
		% Cameron McKenzie, TechTarget.
		% \textit{How to write a Git commit message properly with examples}.\\
		% \url{https://www.theserverside.com/video/Follow-these-git-commit-message-guidelines}.
		% \bibitem{TabellaDelleOreLavorative}
		% Jawa Druids, 2020.
		% \textit{Tabella delle Ore Lavorative}.\\
		% \url{https://docs.google.com/spreadsheets/d/12esX1ISWQOKM-fjuHTLmAzRN0cWltksn7eGiPsBHBI0/edit?usp=sharing}.
		% \bibitem{TrelloJawaDruids}
		% Jawa Druids, 2020.
		% \textit{Trello - dashboard Jawa Druids}.\\
		% \url{https://trello.com/b/hIEOGbE9/jawadruids}.
\end{thebibliography}

%
% \end{document}
